\par\medskip\pdfbookmark[1]{Zu III und IV.}{nexus-note-III-IV}\centerline{Zu III und IV.}\smallskip
Ist $C$ die Gesamtheit aller positiven, nicht eigentlich äquivalenten, eigentlich primitiven Formen mit negativem Determinanten $-D$, und weist man den Variablen jeder dieser Formen zwei zueinander teilerfremde Werte zu, so ist die Anzahl aller Darstellungen einer positiven ganzen Zahl $m$ gleich $\epsilon\psi(m)$. Dabei bezeichnet $\epsilon$ die Anzahl der Lösungen der Gleichung $tt+Duu=1$ und $\psi(m)$ die Anzahl derjenigen Wurzeln $n$ der Kongruenz $nn+D\equiv0\pmod m$, für die die drei Zahlen $m,\ 2n$ und $\frac{nn+D}m$ keinen gemeinsamen Teiler haben (Disqq. Arithm., Art. 180). Für den Faktor gilt $\epsilon$$=4$ für $D=1$, in allen anderen Fällen $=2$. Ist ferner $m=p^\pi p'^{\pi'}p''^{\pi''}\ldots$, wobei $p,p',p''\ldots$ verschiedene Primzahlen sind, so ist $\psi(m)=\psi(p^\pi)\psi(p'^{\pi'})\psi(p''^{\pi''})\ldots$; bezeichnet $\mathfrak A(m)$ die Anzahl der Wurzeln $n$ der Kongruenz $nn+D\equiv0\pmod m$, und verwendet man das von \textsc{Legendre} eingeführte und von \textsc{Jacobi} verallgemeinerte Zeichen, so ist $\psi(p^\pi)=\mathfrak A(p^\pi)=1+\left(\frac{-D}p\right)$, wenn $p$ kein Teiler von $2D$ ist, andernfalls aber $=\mathfrak A(p^\pi-\frac1p\mathfrak A(p^{\pi+1})$; die Anzahl $\mathfrak A(p^\pi)$ lässt sich stets leicht bestimmen (Disqq. Arithm., Art. 104). Für das Folgende genügt jedoch die Bemerkung, dass $\mathfrak A(p^\pi)$ stets von $\pi$ unabhängig wird, sobald $\pi$ eine gewisse Größe überschreitet.

Weist man den Variablen der in der Gesamtheit $C$ enthaltenen Formen ausnahmslos alle ganzzahligen Werte zu (Disqq. Arithm., Art. 181), so wird die Anzahl $(m)$ aller Darstellungen der Zahl $m$ gleich $\epsilon f(m)$, wobei $f(m)=\Sigma\psi\left(\frac m{\mu\mu}\right)$ ist und die Summe über alle quadratischen Teiler $\mu\mu$ der Zahl $m$ zu bilden ist. Daraus folgt unmittelbar
\[
f(m)=f(p^\pi p'^{\pi'}p''^{\pi''}\ldots)=f(p^\pi)f(p'^{\pi'})f(p''^{\pi''})\ldots
\]
und
\[
f(p^\pi)=\psi(p^\pi)+\psi(p^{\pi-2})+\psi(p^{\pi-4})+\ldots
\]
wobei die Reihe so lange fortzusetzen ist, wie die Exponenten $\pi,\ \pi-2,\ \pi-4\ldots$ nicht negativ werden. Wenn $p$ kein Teiler von $2D$ ist, folgt daraus
\[
f(p^\pi)=1+\left(\frac{-D}p\right)+\left(\frac{-D}{pp}\right)+\ldots+\left(\frac{-D}{p^\pi}\right)
\]
und allgemein, wenn $m$ zu $2D$ teilerfremd ist,
\[
f(m)=\Sigma\left(\frac{-D}n\right)
\]
wobei die Summe über alle Teiler $n$ der Zahl $m$ zu bilden ist.

\Needspace{18\baselineskip}
\SourceLeaf{304}
Aus diesen Bemerkungen ergibt sich unmittelbar die Richtigkeit der im Text (III, 1, 2, 3) aufgestellten Sätze über die Anzahl $(m)$, wenn man beim ersten noch die Bedingung hinzufügt, dass $D$ nicht durch $pp$ teilbar sein darf (die Bestimmung der Klassenanzahl ist schon in den Disqq. Arithm., Art. 256, auf den Fall zurückgeführt, in dem $D$ durch kein Quadrat teilbar ist). Zugleich findet man auch ohne diese Einschränkung, dass die unendliche Reihe
\[
1+\frac{f(p)}p+\frac{f(pp)}{pp}+\frac{f(p^3)}{p^3}+\ldots
\]
den Wert
\[
\frac1{1-\frac1p}\quad\text{oder}\quad\frac1{1-.\frac1p}\pdot\frac1{1-\left(\frac{-D}p\right)\frac1p}
\]
hat, je nachdem, ob $2D$ durch die Primzahl $p$ teilbar oder nicht teilbar ist.
\BookRule{.25\textwidth}
